\documentclass[10pt,a4paper]{article}
\usepackage{fontspec,geometry,booktabs,array,multicol,xcolor}
\geometry{margin=17mm}
\setmainfont{DejaVu Serif}
\newfontfamily\arabicfont[Script=Arabic,Scale=1.1]{DejaVu Sans}
\ifdefined\luatexversion
  \newcommand{\ar}[1]{{\arabicfont\textdir TRT #1\textdir TLT}}% LuaLaTeX (untested here)
\else
  \TeXXeTstate=1
  \newcommand{\ar}[1]{{\arabicfont\beginR #1\endR}}% XeLaTeX
\fi
\pagestyle{empty}
\setlength{\parindent}{0pt}
\setlength{\parskip}{5pt}
\begin{document}
{\LARGE English in Arabic Script}\hfill Flyxion\\
{\large Reading and Writing Cheat Sheet}\hfill October 2026

Write English words and sentences from right to left using Arabic letters. Preserve English meaning and word order. These are the current working conventions; some sounds share a spelling, and context helps identify the word.

\begin{multicols}{2}
\textbf{Consonants}\\[3pt]
\begin{tabular}{@{}ll@{}}\toprule
English sound & Letter\\\midrule
b or p & \ar{ب}\\
d & \ar{د}\\
f or v & \ar{ف}\\
hard g & \ar{غ}\\
h & \ar{ه}\\
j (judge) & \ar{ج}\\
k, l, m, n & \ar{ك}\quad\ar{ل}\quad\ar{م}\quad\ar{ن}\\
r, s, t, z & \ar{ر}\quad\ar{س}\quad\ar{ت}\quad\ar{ز}\\
w, y & \ar{و}\quad\ar{ي}\\
sh (she) & \ar{ش}\\
ch (change) & \ar{تْش}\\
th (this) & \ar{ذ}\\
th (think) & \ar{ث}\\
ng (sing) & \ar{نغ}\\
zh (measure) & \ar{ز}\\\bottomrule
\end{tabular}

The assigned English sound can differ from the letter's usual Arabic sound. Use Arabic letters only; no Persian additions. Earlier specimens also used jeem for hard g and ba for v.

\textbf{Vowels and marks}\\[3pt]
\begin{tabular}{@{}lll@{}}\toprule
Cue & Form & Example\\\midrule
short a & \ar{بَ} & \ar{بَتْ} (but)\\
short i/e & \ar{بِ} & \ar{ذِسْ} (this)\\
short u & \ar{بُ} & \ar{بُوكْ} (book)\\
ee & \ar{بِي} & \ar{شِي} (she)\\
oo & \ar{بُو} & \ar{مُونْ} (moon)\\
ay & \ar{بَيْ} & \ar{دَيْ} (day)\\
ai & \ar{بَايْ} & \ar{نَايْتْ} (night)\\
ow & \ar{بَاوْ} & \ar{فْلَاوَرْ} (flower)\\
oy & \ar{بُويْ} & \ar{فُويْسْ} (voice)\\\bottomrule
\end{tabular}

Place a short vowel mark on the preceding consonant. Sukoon \ar{بْ} means no following vowel; the consonant is still pronounced. Long vowel letters normally take no sukoon; glides may take it. Initial vowels use \ar{أَ}, \ar{أُ}, or \ar{إِ}.

\columnbreak
\textbf{Common words}\\[3pt]
\begin{tabular}{@{}ll@{}}\toprule
English & Spelling\\\midrule
the & \ar{ذَا}\\
this & \ar{ذِسْ}\\
that & \ar{ذَتْ}\\
of & \ar{أُفْ}\\
and & \ar{أَنْدْ}\\
in & \ar{إِنْ}\\
to & \ar{تُو}\\
a & \ar{أَ}\\
is & \ar{إِزْ}\\
she & \ar{شِي}\\
it / its & \ar{إِتْ}\quad /\quad\ar{إِتْسْ}\\
as & \ar{أَزْ}\\
into & \ar{إِنْتُو}\\
with & \ar{وِذْ}\\
are & \ar{أَرْ}\\\bottomrule
\end{tabular}

\textbf{Words to recognise}\\[3pt]
\begin{tabular}{@{}ll@{}}\toprule
English & Spelling\\\midrule
book & \ar{بُوكْ}\\
look & \ar{لُوكْ}\\
good & \ar{غُودْ}\\
food & \ar{فُودْ}\\
moon & \ar{مُونْ}\\
foot & \ar{فُوتْ}\\
wood & \ar{وُودْ}\\
green & \ar{غْرِينْ}\\
sleep & \ar{سْلِيبْ}\\
night & \ar{نَايْتْ}\\
voice & \ar{فُويْسْ}\\
change & \ar{تْشَيْنْجْ}\\\bottomrule
\end{tabular}

\textbf{Spelling memory}

Book retains waw because its English spelling has \emph{oo}. The current draft extends this to look, good, foot and wood. These spellings do not distinguish all English vowel qualities. Unstressed vowels can also reflect English spelling.
\end{multicols}
\vspace{2mm}\hrule\vspace{3mm}
\textbf{Read a sentence}\\[4pt]
{\large\ar{ذَا بُوكْ إِزْ رِتِنْ إِنْ إِنْغْلِشْ.}}\\
The book is written in English.

Start with marked words and their English equivalents, then cover the English and read a short passage. Marks may be reduced as recognition develops. This sheet shows working spellings, not a lossless transcription or a keyboard layout.
\end{document}
